\chapter{Magnetostatics and Inductance}
\label{ch:magnetostatics}

\chapterorient{The previous chapter took apart the electrostatic driver and found
a machine of disarmingly few parts: assemble one operator, solve it once per
conductor, and pair the solved family with itself through an energy operator to
read off a capacitance matrix. This chapter builds the magnetostatic driver---the
one that computes \emph{inductance}---and the surprise is how little there is to
build. It is the \emph{same machine}. The skeleton is identical, the solver is the
same conjugate-gradient solver, the reduction is the same Gram verb. Exactly three
things change, and we will name them at the outset and then spend the chapter
earning the right to say they are the \emph{only} three: the partial differential
equation and its field (a vector potential obeying the curl--curl equation, not a
scalar obeying Poisson), the finite-element space that field lives in
($\Hcurl$ edge elements, not $\Hone$ nodal elements), and one scalar weight in the
final reduction (a current normalization, not unity). Everything else---the
once-only assembly, the solve-per-source, the symmetric Gram, even the
positive-definiteness that lets conjugate gradients run---is shared verbatim with
\cref{ch:electrostatics}. By the end you will be able to point at the magnetostatic
pipeline, lay it beside the electrostatic one, and circle precisely the three boxes
that differ.}

\section{The physical question, and the machine we already have}
\label{sec:mag-question}

A handful of current-carrying loops sit in space---traces on a circuit board, turns
of a coil, the signal lines of a connector. Drive loop $i$ with a current $I_i$ and
it produces a magnetic field; that field threads the other loops, and a changing
current in loop $i$ induces a voltage in loop $j$. The bookkeeping of ``how much
flux loop $i$'s current links through loop $j$'' is the \emph{inductance matrix}
$\mat{L}$. Its diagonal $L_{ii}$ is the \emph{self}-inductance of loop $i$; its
off-diagonal $L_{ij}$ is the \emph{mutual} inductance between loops $i$ and $j$.
These are the magnetostatic analogue of the capacitances of
\cref{ch:electrostatics}, and an engineer wants them for the same reason: they are
the lumped circuit parameters that summarize a piece of geometry.

The temptation is to reach for the Biot--Savart law and integrate. We do not. The
whole point of \cref{ch:electrostatics} was that the simulator does not compute a
capacitance by a special-purpose formula; it runs a \emph{generic pipeline} and the
capacitance falls out of a generic reduction. Magnetostatics runs the \emph{same}
pipeline. Before we change anything, let us write down the skeleton we are
inheriting.

\begin{physicsbox}
\textbf{Setup.} $n$ current sources (loops, ports, coil terminals). Drive source
$i$ with unit current $I_i=1$; everything else carries zero current. Solve for the
resulting magnetic vector potential $\vA_i$, from which the magnetic field is
$\vB_i=\Curl\vA_i$. The stored magnetic energy when source $i$ alone is driven is
$U_m^{(i)}=\tfrac{1}{2}L_{ii}I_i^2$, and pairing the field of source $i$ against the
field of source $j$ gives the mutual inductance $L_{ij}$. \emph{We never integrate
Biot--Savart.} We solve a boundary-value problem for $\vA$ and read the inductances
off the solved family---exactly as electrostatics read capacitances off the solved
potential family.
\end{physicsbox}

\Cref{ch:electrostatics} distilled the electrostatic driver to a three-stage
skeleton---the \texttt{assemble}, \texttt{solve}, \texttt{reduce} stages of
\cref{ch:reductions}. Written as a composition, with the fixed-operator solver
corner that \cref{ch:electrostatics} identified:
\begin{lstlisting}[style=l4style]
-- the ELECTROSTATIC driver (Chapter 39), for reference
electrostatic = gram_reduce (w = 1) . solve_family . fe_assemble
  --             ^ reduce: symmetric    ^ one CG solve   ^ assemble
  --               Gram, unit weight       per terminal     ONCE
\end{lstlisting}
Read it right to left, the direction the data flows. \code{fe\_assemble} builds the
operator \emph{once}---the geometry does not change between excitations, so the
matrix is assembled a single time and reused. \code{solve\_family} solves that one
fixed operator against a family of right-hand sides, one per terminal; because the
operator is symmetric positive-definite and never changes, this is the
\emph{fixed-operator} corner of the solve taxonomy (\cref{ch:reductions}), and the
solver is plain conjugate gradients with the operator's preconditioner reused
across the family. \code{gram\_reduce} pairs the solved family with itself through
the energy operator and, with unit weight, returns the capacitance matrix
(\cref{ch:reductions}, eq.~\eqref{eq:cap-gram}).

\begin{keyidea}
The magnetostatic driver is the electrostatic driver with three substitutions and
\emph{nothing else}. The skeleton
\code{reduce} $\circ$ \code{solve\_family} $\circ$ \code{fe\_assemble} is identical;
the assemble-once / solve-per-source / Gram-reduce structure is identical; the
conjugate-gradient solver in the fixed-operator corner is identical. The three
swaps are: \textbf{(1)} the PDE and field---the curl--curl equation
$\Curl(\nu\Curl\vA)=\vJ$ for a vector potential $\vA$, not Poisson for a scalar
$V$; \textbf{(2)} the finite-element space---$\Hcurl$ edge elements, not $\Hone$
nodal elements; \textbf{(3)} the Gram weight---$w=1/(I_iI_j)$, not $w=1$. We spend
the chapter on those three, because they are all there is.
\end{keyidea}

\Cref{fig:mag-pipeline} draws the magnetostatic pipeline laid directly over the
electrostatic one, with the three swapped boxes called out in rust and everything
shared drawn in the muted shared palette. It is the hero figure of the chapter:
keep it in view as we take the swaps one at a time.

\begin{figure}[t]
\centering
\resizebox{\linewidth}{!}{%
\begin{tikzpicture}[node distance=6mm and 13mm, font=\footnotesize]
  % --- shared config input ---
  \node[data, align=center, text width=20mm] (cfg) {config\\\footnotesize geometry,
    sources};
  % --- assemble (SWAPPED: operator) ---
  \node[stage, fill=simBgRust, draw=simRust, right=of cfg, align=center,
    text width=24mm] (asm) {\textbf{fe\_assemble}\\\footnotesize \textsc{swap 1+2}};
  % --- solve (SHARED) ---
  \node[stage, fill=simBgGray, draw=simGray, right=of asm, align=center,
    text width=24mm] (slv) {\textbf{solve\_family}\\\footnotesize CG, fixed op
    \emph{(shared)}};
  % --- reduce (SWAPPED: weight) ---
  \node[stage, fill=simBgRust, draw=simRust, right=of slv, align=center,
    text width=24mm] (red) {\textbf{gram\_reduce}\\\footnotesize \textsc{swap 3}};
  % --- product ---
  \node[product, right=of red, align=center, text width=18mm] (prod)
    {$\mat{L}$\\\footnotesize inductance};
  \draw[flow] (cfg) -- (asm);
  \draw[flow] (asm) -- (slv);
  \draw[flow] (slv) -- (red);
  \draw[flow] (red) -- (prod);
  % --- swap annotations underneath ---
  \node[font=\scriptsize, simRust, align=center, below=8mm of asm]
    (s12) {\textsc{swap 1}: curl--curl PDE\\\textsc{swap 2}: $\Hcurl$ space\\
      $\Curl(\nu\Curl\vA)=\vJ$};
  \node[font=\scriptsize, simGray, align=center, below=8mm of slv]
    (ss) {identical to\\\cref{ch:electrostatics}:\\SPD operator, CG};
  \node[font=\scriptsize, simRust, align=center, below=8mm of red]
    (s3) {\textsc{swap 3}: weight\\$w=\dfrac{1}{I_iI_j}$\\(was $w=1$)};
  \draw[refedge] (asm) -- (s12);
  \draw[refedge] (slv) -- (ss);
  \draw[refedge] (red) -- (s3);
\end{tikzpicture}%
}
\caption[The magnetostatic pipeline, three swaps from electrostatics]{The
magnetostatic driver drawn over the electrostatic skeleton of
\cref{ch:electrostatics}. The three rust boxes are the \emph{only} things that
change. \textsc{Swap 1+2} live in \code{fe\_assemble}: the partial differential
equation becomes the curl--curl equation $\Curl(\nu\Curl\vA)=\vJ$ (\textsc{swap 1})
and the finite-element space becomes $\Hcurl$ edge elements (\textsc{swap 2}), so a
different stiffness operator is assembled. \textsc{Swap 3} lives in
\code{gram\_reduce}: the Gram weight becomes the current normalization
$w=1/(I_iI_j)$. The grey \code{solve\_family} box is byte-for-byte the electrostatic
solve---the same conjugate-gradient solver on a symmetric positive-definite,
fixed-across-the-family operator. Assemble once, solve per source, reduce by
pairing: the machine is shared; only the three rust boxes are magnetostatic.}
\label{fig:mag-pipeline}
\end{figure}

\section{Swap 1: the curl--curl equation and its field}
\label{sec:swap1}

Electrostatics solves Poisson's equation for a scalar potential. Magnetostatics
solves the \emph{curl--curl equation} for a vector potential. This is the first
swap, and it is forced by the physics: a magnetic field is a fundamentally
different kind of object from an electric one, and the de Rham complex
(\cref{ch:derham}) told us exactly which potential each admits.

\subsection{From Maxwell to curl--curl}

Two of Maxwell's laws govern magnetostatics. Amp\`ere's law (in the static,
source-driven regime) relates the magnetic field $\vH$ to the current density
$\vJ$, and the constitutive law relates $\vH$ to $\vB$ through the permeability:
\[
  \Curl\vH = \vJ, \qquad \vB = \mu\,\vH, \qquad \Div\vB = 0 .
\]
The last equation, $\Div\vB=0$---no magnetic monopoles---is the structural fact.
By exactness of the de Rham complex at the $\Hdiv$ rung (\cref{ch:derham},
\cref{thm:exact}), a divergence-free field is a curl: there exists a
\emph{magnetic vector potential} $\vA$ with
\[
  \boxed{\;\vB = \Curl\vA\;}
\]
Substitute this into Amp\`ere's law. Writing $\nu=1/\mu$ for the
\emph{reluctivity} (the magnetic analogue of conductivity---``how much a material
resists magnetic flux''), we have $\vH=\nu\vB=\nu\Curl\vA$, and Amp\`ere's law
becomes
\begin{equation}
  \boxed{\;\Curl\!\left(\nu\,\Curl\vA\right) = \vJ\;}
  \label{eq:curlcurl}
\end{equation}
the \emph{curl--curl equation} for the vector potential $\vA$, driven by the
prescribed source current $\vJ$. This is the magnetostatic governing equation
(\cref{ch:reductions-pde}). \Cref{fig:loop-field} shows the physical picture: a
current loop, the vector potential $\vA$ that circulates with the current, and the
field $\vB=\Curl\vA$ that loops through.

\begin{figure}[t]
\centering
\begin{tikzpicture}[scale=1.0, font=\footnotesize]
  % --- the current loop (an ellipse seen at an angle) ---
  \draw[simRust, very thick] (0,0) ellipse (1.5 and 0.55);
  % current direction arrows ON the loop
  \draw[-{Stealth[length=2mm]}, simRust, very thick] (1.5,0.05)--(1.5,-0.05);
  \draw[-{Stealth[length=2mm]}, simRust, very thick] (-1.5,-0.05)--(-1.5,0.05);
  \node[simRust] at (0,-0.85) {current loop, $I$};
  \node[simRust, font=\scriptsize] at (1.95,0.0) {$\vJ$};
  % --- B field lines threading the loop (vertical-ish loops) ---
  \draw[simBlue, opacity=0.4,  thick] (0,0) ellipse [x radius=0.45, y radius=0.675];
  \draw[simBlue, opacity=0.55, thick] (0,0) ellipse [x radius=0.85, y radius=1.275];
  \draw[simBlue, opacity=0.4,  thick] (0,0) ellipse [x radius=1.25, y radius=1.875];
  % B arrow up through the center
  \draw[-{Stealth[length=2.4mm]}, simBlue, thick] (0,-1.7)--(0,1.7);
  \node[simBlue] at (0.35,1.55) {$\vB$};
  \node[simBlue, font=\scriptsize, align=center] at (2.4,1.3)
    {$\vB=\Curl\vA$\\(field lines\\thread the loop)};
  % --- A field circulating WITH the current (faint, around the wire) ---
  \node[simGreen, font=\scriptsize, align=center] at (-2.5,0.9)
    {$\vA$ circulates\\\emph{with} the current\\(same direction as $\vJ$)};
  \draw[-{Stealth[length=1.8mm]}, simGreen, thick] (-1.95,0.35) arc (160:380:0.32 and 0.20);
\end{tikzpicture}
\caption[A current loop and its potential]{The physical picture behind the
curl--curl equation. A current $I$ runs around the loop (rust). The magnetic vector
potential $\vA$ (green) circulates \emph{with} the current---it lives one rung up
the de Rham complex from $\vB$, in $\Hcurl$, the space of circulations
(\cref{ch:derham}). Its curl is the magnetic field $\vB=\Curl\vA$ (blue), whose
field lines thread \emph{through} the loop. The simulator solves
$\Curl(\nu\Curl\vA)=\vJ$ for $\vA$ on the whole domain, then post-processes
$\vB=\Curl\vA$. Compare the electrostatic picture, where the scalar potential $V$
sits at \emph{vertices} and $\vE=-\Grad V$.}
\label{fig:loop-field}
\end{figure}

\subsection{The contrast with Poisson, term by term}

Lay \cref{eq:curlcurl} beside the electrostatic Poisson equation
$-\Div(\varepsilon\,\Grad V)=\rho$ of \cref{ch:electrostatics}. The two are the
\emph{same shape}---a second-order, self-adjoint, source-driven elliptic
equation---built from a different rung of the de Rham complex:
\[
  \underbrace{-\Div(\varepsilon\,\Grad V)=\rho}_{\text{electrostatics}}
  \qquad\longleftrightarrow\qquad
  \underbrace{\Curl(\nu\,\Curl\vA)=\vJ}_{\text{magnetostatics}} .
\]
The gradient is replaced by the curl, the divergence (its adjoint) by the curl
again (the curl is its own adjoint up to boundary terms), the scalar coefficient
$\varepsilon$ by the scalar coefficient $\nu$, the scalar unknown $V$ by the vector
unknown $\vA$, and the scalar source $\rho$ by the vector source $\vJ$. The weak
form (\cref{ch:weak}) follows the identical recipe: multiply by a test field,
integrate by parts once, and move one curl onto the test field. The bilinear form
is
\begin{equation}
  a(\vA,\vF) = \int_\dom \nu\,(\Curl\vA)\cdot(\Curl\vF)\dV,
  \qquad
  \ell(\vF) = \int_\dom \vJ\cdot\vF \dV,
  \label{eq:weak-curlcurl}
\end{equation}
to be compared with the electrostatic
$a(V,W)=\int\varepsilon\,\Grad V\cdot\Grad W$. Discretizing
\cref{eq:weak-curlcurl} by Galerkin projection (\cref{ch:galerkin}) produces the
\emph{curl--curl stiffness matrix} $\mat{K}$, the magnetostatic energy operator:
\begin{equation}
  \mat{K}_{ab} = \int_\dom \nu\,(\Curl\boldsymbol\phi_a)\cdot(\Curl\boldsymbol\phi_b)\dV,
  \label{eq:stiff-curlcurl}
\end{equation}
where $\boldsymbol\phi_a$ are the basis functions of the finite-element space---and
\emph{which} space those are is the second swap.

\begin{notationbox}
The same letter $\mat{K}$ denotes the energy/stiffness operator in both drivers,
because it plays the identical structural role: it is the symmetric
positive-(semi)definite operator that \code{solve\_family} inverts and that
\code{gram\_reduce} pairs the family through. In electrostatics
$\mat{K}_{ab}=\int\varepsilon\,\Grad\phi_a\cdot\Grad\phi_b$; in magnetostatics
$\mat{K}_{ab}=\int\nu\,\Curl\boldsymbol\phi_a\cdot\Curl\boldsymbol\phi_b$. One
symbol, one role, two integrands---the swap is entirely inside the integrand.
\end{notationbox}

\section{Swap 2: the field lives in \texorpdfstring{$\Hcurl$}{H(curl)}, on edge elements}
\label{sec:swap2}

The unknown is now a vector field $\vA$, and the single most consequential decision
in the whole chapter is which finite-element space to put it in. The naive choice
is to take three copies of the nodal $\Hone$ space---one per Cartesian component---
and store $\vA$'s components at vertices, as we stored $V$. \emph{This is wrong},
and the way it is wrong is the deepest lesson of \cref{ch:functionspaces} and
\cref{ch:derham}. The correct space is $\Hcurl$, realized by \emph{edge elements}
(N\'ed\'elec elements). This is the second swap.

\subsection{Why a vector potential wants edge elements}

The de Rham complex (\cref{ch:derham}) names the home of every electromagnetic
field, and it is unambiguous: $\vA$ generates $\vB$ by $\vB=\Curl\vA$, so $\vA$ sits
at the $\Hcurl$ rung---the space whose degrees of freedom are \emph{edge
circulations} $\int_e\vA\cdot d\vec\ell$, and whose defining continuity is that the
\emph{tangential component} of $\vA$ is continuous across element faces while the
normal component may jump. That tangential continuity is exactly what a vector
potential physically demands. The tangential component of $\vA$ carries its
circulation, hence its curl, hence the magnetic field $\vB$; the normal component
carries the gauge (\cref{ch:derham}, \cref{sec:potentials}) and is physically
free to jump. \Cref{fig:nodal-vs-edge} draws the two choices side by side.

\begin{figure}[t]
\centering
\begin{tikzpicture}[scale=1.0, font=\scriptsize]
% ===== LEFT: nodal vector (WRONG) =====
\begin{scope}
  \node[font=\footnotesize\bfseries, simRust] at (1.2,2.35) {nodal vector (wrong)};
  \coordinate (a) at (0,0); \coordinate (b) at (2.4,0); \coordinate (c) at (1.2,1.9);
  \fill[simBgRust] (a)--(b)--(c)--cycle;
  \draw[simInk, thick] (a)--(b)--(c)--cycle;
  % full vector at each vertex (all 3 components stored)
  \foreach \p in {a,b,c}{
    \fill[simRust] (\p) circle (2.6pt);
    \draw[-{Stealth[length=1.6mm]}, simRust, thick] (\p)++(0,0)--++(0.45,0.28);
    \draw[-{Stealth[length=1.6mm]}, simRust, thick] (\p)++(0,0)--++(-0.10,0.50);
  }
  \node[simRust, align=center] at (1.2,-0.55) {\emph{both} components\\at each vertex};
  \node[align=center, text=simInk] at (1.2,-1.35)
    {forces \textbf{full} continuity\\$\Rightarrow$ spurious gradient modes};
\end{scope}
% ===== RIGHT: edge (RIGHT) =====
\begin{scope}[shift={(5.6,0)}]
  \node[font=\footnotesize\bfseries, simGreen] at (1.2,2.35) {edge / N\'ed\'elec (right)};
  \coordinate (a) at (0,0); \coordinate (b) at (2.4,0); \coordinate (c) at (1.2,1.9);
  \fill[simBgGreen] (a)--(b)--(c)--cycle;
  \draw[simInk, thick] (a)--(b)--(c)--cycle;
  % one tangential circulation per edge
  \draw[-{Stealth[length=1.8mm]}, simGreen, very thick] ($(a)!0.32!(b)$)--($(a)!0.68!(b)$);
  \draw[-{Stealth[length=1.8mm]}, simGreen, very thick] ($(b)!0.32!(c)$)--($(b)!0.68!(c)$);
  \draw[-{Stealth[length=1.8mm]}, simGreen, very thick] ($(c)!0.32!(a)$)--($(c)!0.68!(a)$);
  \node[simGreen, align=center] at (1.2,-0.55) {\emph{one tangential}\\DoF per edge};
  \node[align=center, text=simInk] at (1.2,-1.35)
    {tangential continuity only\\$\Rightarrow$ exact discrete complex};
\end{scope}
\end{tikzpicture}
\caption[Nodal vector versus edge elements]{Two ways to discretize a vector field,
and why only one works. \emph{Left:} the naive nodal-vector element stores the full
vector (all components) at each \emph{vertex}, forcing \emph{full} continuity of
$\vA$ across faces---more continuity than the physics allows. The excess continuity
admits non-physical fields and (as the next section shows) pollutes the curl--curl
operator's null space with spurious modes. \emph{Right:} the edge (N\'ed\'elec)
element of $\Hcurl$ stores one \emph{tangential circulation} per \emph{edge},
imposing only tangential continuity---exactly what a vector potential requires. The
edge space sits in an \emph{exact} discrete de Rham complex (\cref{ch:derham},
\cref{thm:discrete-exact}), so its discrete curl has precisely the kernel the
physics demands, with no spurious extras. This is \emph{the} reason edge elements
exist (\cref{ch:bc}).}
\label{fig:nodal-vs-edge}
\end{figure}

The nodal-vector choice imposes \emph{full} continuity of $\vA$---both tangential
and normal components glued across every face. That is strictly more continuity
than $\vB=\Curl\vA$ requires, and the excess is not harmless. A field
discretization that is ``too continuous'' for its operator admits non-physical
solutions; in the curl--curl setting these surface as \emph{spurious modes}
(\cref{ch:functionspaces}, \cref{ch:bc})---finite-element fields with no physical
counterpart that nonetheless satisfy the discrete equations and contaminate the
spectrum. \Cref{ch:derham} diagnosed the cause precisely: nodal-vector elements do
not sit in an exact discrete de Rham complex, so their discrete curl has the
\emph{wrong} kernel, and the wrong kernel is the breeding ground of spurious modes.
Edge elements sit in an exact discrete complex (\cref{ch:derham},
\cref{thm:discrete-exact}); their discrete curl's kernel is \emph{exactly} the range
of the discrete gradient---the gauge fields, and nothing else.

\begin{pitfall}
Do \emph{not} discretize the magnetic vector potential with three copies of nodal
$\Hone$ elements ``one per component.'' It looks natural---it is how you would store
a vector in a general-purpose array---but it imposes more continuity than $\vA$
admits and breaks the exact discrete de Rham complex. The cost is spurious modes
(\cref{ch:bc}): non-physical fields polluting the curl--curl operator's null space,
which corrupt eigenvalue computations and stall iterative solvers. The vector
potential lives in $\Hcurl$, full stop, and $\Hcurl$ is realized by edge elements.
This is not an efficiency preference; it is a correctness requirement.
\end{pitfall}

\subsection{The curl--curl operator has a large null space}
\label{sec:nullspace}

Choosing edge elements fixes \emph{which} kernel the operator has, but it does not
make the kernel small. The curl--curl operator carries an enormous null space, and
acknowledging it is the price of working with vector potentials. The reason is the
gauge freedom of \cref{ch:derham}: for any scalar field $\psi$,
\[
  \Curl\!\left(\nu\,\Curl(\Grad\psi)\right) = \Curl\!\left(\nu\cdot\mathbf 0\right)
  = \mathbf 0,
  \qquad\text{since } \Curl\Grad\psi = \mathbf 0 .
\]
Every gradient field is annihilated by the curl--curl operator. Discretely, every
discrete-gradient vector $\mat{G}\psi_h$ lies in the null space of the curl--curl
stiffness $\mat{K}$ (\cref{eq:stiff-curlcurl}). So $\mat{K}$ is only \emph{positive
semi-definite}: it is singular, with a null space that is the entire range of the
discrete gradient $\mat{G}$ (\cref{ch:derham}, \cref{def:discrete-grad}). This is
the same gradient null space the eigenmode solver must project out
(\cref{ch:bc}) and the same null space a plain smoother cannot damp
(\cref{ch:smoothers}). \Cref{fig:curlcurl-null} draws it.

\begin{figure}[t]
\centering
\begin{tikzpicture}[font=\footnotesize]
  % the edge-dof space as a big rounded box
  \node[draw=simInk, thick, rounded corners=3pt, minimum width=66mm,
    minimum height=30mm, fill=simBgGray] (box) at (0,0) {};
  \node[font=\scriptsize\itshape, simInk, anchor=north west] at (box.north west)
    {\;edge-DoF space $\Hcurl_h$};
  % the gradient null space as an inner blob
  \node[draw=simBlue, thick, rounded corners=3pt, fill=simBgBlue,
    minimum width=30mm, minimum height=20mm, align=center] (null) at (-1.4,-0.3)
    {$\ker(\mat{K})$\\$=\operatorname{range}(\mat{G})$\\[1pt]
     \scriptsize gradient fields $\Grad\psi$\\\scriptsize ($\mat{K}\,\Grad\psi=0$)};
  % the physical / div-free complement
  \node[font=\scriptsize, simGreen, align=center, anchor=west] at (0.7,0.55)
    {physical part\\(divergence-free)\\$\mat{K}$ is SPD \emph{here}};
  \node[font=\scriptsize, simRust, align=center, below=1mm of box.south]
    {curl--curl $\mat{K}$ is only \emph{positive semi-definite}: singular on
     $\operatorname{range}(\mat{G})$};
\end{tikzpicture}
\caption[The curl--curl gradient null space]{The curl--curl stiffness operator
$\mat{K}$ on the edge-DoF space. Its null space (blue) is the \emph{entire} range of
the discrete gradient $\mat{G}$---every gradient field $\Grad\psi$ has zero curl, so
$\mat{K}\,\Grad\psi=0$ (the gauge freedom of \cref{ch:derham}). The operator is
therefore only positive \emph{semi}-definite, singular on a large subspace, and a
plain conjugate-gradient solve does not automatically converge. Two cures
(\cref{sec:solvecare}): make the right-hand side $\vJ$ \emph{consistent}
(divergence-free, so it has no component along the null space, and CG runs in the
SPD complement), or precondition with an auxiliary-space method (AMS,
\cref{ch:smoothers}) that handles the gradient block explicitly. Either way the
\emph{solver itself} is the shared CG of \cref{ch:electrostatics}; only its setup
accounts for the null space.}
\label{fig:curlcurl-null}
\end{figure}

\subsection{What this costs the solve---and what it does not}
\label{sec:solvecare}

A singular operator threatens \cref{fig:mag-pipeline}'s claim that the solve is
``identical to electrostatics.'' It is worth being precise about what changes. The
electrostatic $\mat{K}$ is positive \emph{definite} (after grounding a reference
node), so conjugate gradients converges with no further thought. The magnetostatic
$\mat{K}$ is positive \emph{semi}-definite, so a naive CG could wander in the null
space. Two standard cures restore convergence, and crucially \emph{neither changes
the solver}:

\begin{itemize}[nosep]
\item \textbf{Consistent source.} If the driving current $\vJ$ is
  divergence-free---$\Div\vJ=0$, the physical condition for a steady current that
  neither accumulates nor depletes charge---then $\vJ$ has \emph{no component along
  the null space} (the null space is gradients, and a divergence-free field is
  orthogonal to gradients). Conjugate gradients then runs entirely in the
  positive-definite complement, never excited along the singular directions, and
  converges to the (gauge-fixed) solution exactly as in electrostatics. This is the
  cheapest cure and the one Palace prefers: enforce a discretely-consistent source
  and the existing CG just works.
\item \textbf{Auxiliary-space preconditioning (AMS).} When robustness across mesh
  and coefficient is needed, the curl--curl solve is preconditioned by the
  \emph{auxiliary-space Maxwell} (AMS) preconditioner~\citep{hiptmair1998}, which
  splits the error into its divergence-free part and its gradient part and
  preconditions each on the appropriate auxiliary space (the nodal $\Hone$ space for
  the gradient block, via the discrete gradient $\mat{G}$). AMS is exactly the
  multigrid-for-edge-elements machinery of \cref{ch:smoothers}: it exists
  \emph{because} a plain smoother cannot see the gradient null space.
\end{itemize}

\begin{keyidea}
The curl--curl operator's gradient null space changes the solve's \emph{setup}, not
its \emph{solver}. Make the source consistent ($\Div\vJ=0$) so the null space is
never excited, or precondition with AMS (\cref{ch:smoothers}) so the gradient block
is handled explicitly---but in either case the iteration is the same
conjugate-gradient method on a symmetric, fixed-across-the-family operator that
electrostatics used. \code{solve\_family} in \cref{fig:mag-pipeline} is genuinely
shared; the null space is accounted for in how $\mat{K}$ and $\vJ$ are prepared, not
in a new algorithm.
\end{keyidea}

\section{Swap 3: the Gram weight is a current normalization}
\label{sec:swap3}

The solved family $\{\vA_0,\dots,\vA_{n-1}\}$ in hand, the reduce stage runs the
\emph{same} verb electrostatics ran---\code{gram\_reduce} of \cref{ch:reductions}---
and the third and final swap is a single scalar. This is the punchline of the
chapter, and it is the cleanest instance in the book of ``one verb, two physical
quantities.''

\subsection{One verb, both quantities}

Recall the Gram reduction (\cref{ch:reductions}, eq.~\eqref{eq:gram-entry}). Given a
solution family, the energy operator $\mat{K}$, and a scalar weight $w(i,j)$, it
computes the symmetric matrix $G_{ij}=w(i,j)\,\vx_j^{\!\top}\mat{K}\,\vx_i$:
\begin{lstlisting}[style=l4style]
gram_reduce :: LinOp -> [Tensor] -> (Int -> Int -> Scalar) -> Matrix
gram_reduce k xs w =
  symmetric_from_upper                              -- mirror lower triangle
    [ [ w i j * (xs!!j `bilinear` k $ xs!!i) | j <- [i .. m-1] ]
      | i <- [0 .. m-1] ]                           -- map over upper-triangle pairs
  where m = length xs
\end{lstlisting}
This is verbatim the reduction electrostatics used. The energy operator $\mat{K}$,
the symmetric pairing $\vx_j^{\!\top}\mat{K}\vx_i$, the upper-triangle map, the
mirror---all shared. Capacitance plugged in $w(i,j)=1$ and the per-terminal
potential family, getting $C_{ij}=V_j^{\!\top}\mat{K}\,V_i$
(\cref{ch:reductions}, eq.~\eqref{eq:cap-gram}). Inductance plugs in the
\emph{current-normalized} weight and the per-source vector-potential family:
\begin{equation}
  \boxed{\;w(i,j) = \frac{1}{I_i\,I_j},
  \qquad
  L_{ij} = \frac{1}{I_i I_j}\,\vA_j^{\!\top}\,\mat{K}\,\vA_i\;}
  \label{eq:ind-gram-local}
\end{equation}
with $\mat{K}$ the $\nu$-weighted curl--curl energy operator of
\cref{eq:stiff-curlcurl}. \Cref{fig:one-verb} sets the two side by side.

\begin{figure}[t]
\centering
\begin{tikzpicture}[font=\footnotesize]
% --- shared verb in the center ---
\node[stage, fill=simBgGold, draw=simGold, align=center, text width=34mm]
  (verb) at (0,0) {\textbf{gram\_reduce}\\[1pt]
  $G_{ij}=w(i,j)\,\vx_j^{\!\top}\mat{K}\,\vx_i$\\[1pt]\footnotesize
  symmetric, upper-triangle map};
% --- left branch: capacitance ---
\node[product, align=center, text width=30mm, above left=11mm and 14mm of verb]
  (cap) {\textbf{capacitance} $\mat{C}$\\[1pt]\footnotesize
  $w=1$\\family: potentials $V_i\in\Hone$};
% --- right branch: inductance ---
\node[product, fill=simBgBlue, draw=simBlue, align=center, text width=30mm,
  above right=11mm and 14mm of verb]
  (ind) {\textbf{inductance} $\mat{L}$\\[1pt]\footnotesize
  $w=\dfrac{1}{I_iI_j}$\\family: potentials $\vA_i\in\Hcurl$};
\draw[flow] (verb) -- node[left, font=\scriptsize, simGray]{$w\!=\!1$} (cap);
\draw[flow] (verb) -- node[right, font=\scriptsize, simBlue]
  {$w\!=\!\tfrac{1}{I_iI_j}$} (ind);
% --- the "everything else identical" band below ---
\node[font=\scriptsize\itshape, simGray, align=center, below=7mm of verb,
  text width=72mm]
  {identical: the operator-weighted bilinear form $\vx_j^{\!\top}\mat{K}\vx_i$, the
   structural symmetry $G_{ij}=G_{ji}$, the embarrassingly parallel map. \emph{The
   only differences are the scalar weight $w$ and which family is fed in.}};
\end{tikzpicture}
\caption[One verb, two physical quantities]{The single Gram verb
\code{gram\_reduce} (\cref{ch:reductions}) produces \emph{both} capacitance and
inductance. Feed it the per-terminal potential family with unit weight $w=1$ and it
returns the capacitance matrix $\mat{C}$ (left); feed it the per-source
vector-potential family with the current-normalized weight $w=1/(I_iI_j)$ and it
returns the inductance matrix $\mat{L}$ (right). The operator-weighted bilinear
form, the structural symmetry, and the parallel map are byte-for-byte identical---
the two physical quantities differ \emph{only} in the scalar weight and in which
solved family is supplied. This is \textsc{swap 3}, and it is the entire difference
in the reduce stage.}
\label{fig:one-verb}
\end{figure}

\subsection{Why inductance needs the normalization and capacitance does not}

The asymmetry is real but shallow, and it is worth saying plainly. In
electrostatics the excitation is a \emph{voltage}, and we choose it to be unity
($V_i=1$); the energy bilinear form $V_j^{\!\top}\mat{K}V_i$ is already
per-unit-voltage, so dividing by $V_iV_j=1$ does nothing and the weight is $1$. In
magnetostatics the natural excitation is a \emph{current} $I_i$, and the magnetic
energy bilinear form $\vA_j^{\!\top}\mat{K}\vA_i$ scales with the \emph{product of
the driving currents} (double the current, double $\vA$, quadruple the energy). To
recover the geometry-only inductance $L=2U_m/I^2$ (\cref{ch:maxwell}), we must
divide out that current scaling: $w=1/(I_iI_j)$. The weight is precisely the
``per-unit-excitation'' rescaling---invisible when the excitation is chosen to be
unity, explicit when it is a physical current we do not get to set to $1$. Same
energy form, same machine; one quantity happens to normalize trivially and the other
does not.

\begin{keyidea}
Capacitance and inductance are the \emph{same Gram reduction} differing in one
scalar weight: $w=1$ for capacitance, $w=1/(I_iI_j)$ for inductance. The diagonal
$L_{ii}=\vA_i^{\!\top}\mat{K}\vA_i/I_i^2$ is twice the stored magnetic energy per
unit current squared---the self-inductance $L=2U_m/I^2$. The off-diagonal $L_{ij}$
is the mutual inductance, the flux of source $j$'s field linked by loop $i$,
normalized by both currents. The current normalization is the \emph{only} thing the
reduce stage does differently from capacitance.
\end{keyidea}

\section{The composition, and reading the result}
\label{sec:compose}

We can now write the magnetostatic driver as a single composition and put it beside
the electrostatic one. They are the same expression with three substitutions:

\begin{lstlisting}[style=l4style]
-- ELECTROSTATIC (Chapter 39):
electrostatic = gram_reduce (w = 1)          . solve_family . fe_assemble{- Poisson, H^1 -}

-- MAGNETOSTATIC (this chapter), three swaps marked:
magnetostatic = gram_reduce (w = \i j -> 1/(I i * I j))   -- SWAP 3: current weight
              . solve_family                              --   (shared CG, fixed op)
              . fe_assemble{- curl-curl, H(curl) -}       -- SWAP 1+2: PDE + space

-- read B from the solved A:
postprocess A_i = curl A_i      -- B_i = Curl A_i  (one rung down the complex)
\end{lstlisting}
The structure is identical; the differences are exactly the curl--curl-on-$\Hcurl$
assembly (swaps 1 and 2, both inside \code{fe\_assemble}) and the current-normalized
Gram weight (swap 3, inside \code{gram\_reduce}). The forward reference to the
eigenmode driver (\cref{ch:eigenmodes}) is worth flagging now: it reuses this very
same curl--curl $\mat{K}$ on this same $\Hcurl$ space---the magnetostatic assembly
is the eigenmode assembly---and the gradient null space we met in
\cref{sec:nullspace} is precisely what the eigenmode solver must project out.

\subsection{The inductance matrix: self and mutual}

The product is the symmetric $n\times n$ inductance matrix $\mat{L}$. Reading it:

\begin{itemize}[nosep]
\item \textbf{Diagonal $L_{ii}$ (self-inductance).} The energy stored when source
  $i$ alone carries unit current, doubled: $L_{ii}=\vA_i^{\!\top}\mat{K}\vA_i/I_i^2
  =2U_m^{(i)}/I_i^2$. Always positive (the energy operator is positive
  semi-definite and $\vA_i$ is in its definite complement)---self-inductance is
  never negative.
\item \textbf{Off-diagonal $L_{ij}$ (mutual inductance).} The flux of source $j$'s
  field linked by loop $i$, per unit of each current:
  $L_{ij}=\vA_j^{\!\top}\mat{K}\vA_i/(I_iI_j)$. Its \emph{sign} encodes the relative
  orientation of the two loops---positive when their fields reinforce, negative when
  they oppose. Symmetric: $L_{ij}=L_{ji}$, the reciprocity that the Gram's structural
  symmetry guarantees for free.
\end{itemize}

\Cref{fig:Lmatrix} shows the matrix and its physical reading.

\begin{figure}[t]
\centering
\begin{tikzpicture}[font=\footnotesize]
  % --- the 3x3 inductance matrix with labeled cells ---
  \foreach \i in {0,1,2}{
    \foreach \j in {0,1,2}{
      \node[draw=simGray, semithick,
        fill={\ifnum\i=\j simBgBlue\else simBgGold\fi},
        minimum size=10mm, inner sep=0pt] (L\i\j) at (\j*1.0, -\i*1.0) {};}}
  \node at (0,0)    {$L_{11}$};
  \node at (1,-1)   {$L_{22}$};
  \node at (2,-2)   {$L_{33}$};
  \node at (1,0)    {$L_{12}$};
  \node at (2,0)    {$L_{13}$};
  \node at (0,-1)   {$L_{21}$};
  \node at (2,-1)   {$L_{23}$};
  \node at (0,-2)   {$L_{31}$};
  \node at (1,-2)   {$L_{32}$};
  % symmetry mirror
  \draw[densely dashed, simRust, semithick] (-0.55,0.55)--(2.55,-2.55);
  % annotations
  \node[simBlue, align=left, anchor=west, font=\scriptsize] at (3.0,-0.1)
    {\textbf{diagonal}: self-inductance\\$L_{ii}=2U_m^{(i)}/I_i^2>0$};
  \node[simGold, align=left, anchor=west, font=\scriptsize] at (3.0,-1.4)
    {\textbf{off-diagonal}: mutual\\$L_{ij}=\vA_j^{\!\top}\mat{K}\vA_i/(I_iI_j)$\\
     sign $=$ relative orientation};
  \node[simRust, align=left, anchor=west, font=\scriptsize] at (3.0,-2.4)
    {\textbf{symmetric}: $L_{ij}=L_{ji}$\\(reciprocity, free from the Gram)};
\end{tikzpicture}
\caption[The inductance matrix]{The inductance matrix $\mat{L}$, the magnetostatic
product. The diagonal entries (blue) are the self-inductances $L_{ii}=2U_m^{(i)}/
I_i^2$, always positive. The off-diagonal entries (gold) are the mutual inductances
$L_{ij}$, whose sign records whether loops $i$ and $j$'s fields reinforce
($+$) or oppose ($-$). The matrix is symmetric (dashed mirror): $L_{ij}=L_{ji}$, a
reciprocity that the Gram reduction's structural symmetry supplies at no cost---only
the upper triangle is computed (\cref{ch:reductions}). This is the same matrix
\emph{shape} as the capacitance matrix of \cref{ch:electrostatics}; the entries are
inductances because the family is vector potentials and the weight is
current-normalized.}
\label{fig:Lmatrix}
\end{figure}

\begin{remark}[Loop inductance, partial inductance, and what we measured]
We have computed the \emph{loop} inductance of a set of complete current loops: each
$\vA_i$ is the field of a full, closed, divergence-free current path, so $L_{ij}$ is
a well-defined, gauge-independent number. A related notion, \emph{partial}
inductance, ascribes inductance to a single open segment of wire (useful for
circuit extraction from layout); it is not gauge-independent in isolation and only
the loop sums it builds are physical. The magnetostatic driver as composed here
returns loop inductances directly. The partial-inductance bookkeeping is a
reduce-stage variation---a different weighting and pairing of the same solved
family---and we note it only to flag that ``inductance matrix'' can mean slightly
different things downstream; the field solve underneath is the one in this chapter.
\end{remark}

\section{Worked examples}

\subsection{A two-loop inductance, end to end}

We trace the magnetostatic pipeline through a two-loop problem, in deliberate
parallel with the two-terminal \emph{capacitance} computed in \cref{ch:reductions}
(\cref{wex:cap2x2}). Same machine, same numbers shape---inductances out instead of
capacitances, because of the three swaps.

\begin{workedexample}{A $2\times2$ inductance from a current-normalized Gram}{ind2x2}
\textbf{Setup.} Two current loops, driven one at a time at unit current
$I_1=I_2=1\,\si{\ampere}$. The magnetostatic driver assembles the curl--curl energy
operator $\mat{K}$ \emph{once} (swaps 1+2: curl--curl on $\Hcurl$), then solves it
twice (\code{solve\_family}), once per loop, for the vector-potential fields $\vA_1$
and $\vA_2$. We are handed the three operator-weighted pairings that the solve
produced:
\[
  \vA_1^{\!\top}\mat{K}\,\vA_1 = 8,\qquad
  \vA_2^{\!\top}\mat{K}\,\vA_2 = 6,\qquad
  \vA_2^{\!\top}\mat{K}\,\vA_1 = 2
\]
(in nanohenry-equivalent units; these are the magnetic-energy bilinear forms, twice
the stored energy in each configuration).

\textbf{The Gram reduction (swap 3).} The inductance is \code{gram\_reduce} with the
current-normalized weight $w(i,j)=1/(I_iI_j)$ (\cref{eq:ind-gram-local}). With unit
currents, $w\equiv 1$ \emph{numerically}---but the normalization is structurally
present, and would rescale every entry if the currents were not $1$. The verb
computes the upper triangle and mirrors:
\[
  L_{11} = \frac{\vA_1^{\!\top}\mat{K}\vA_1}{I_1^2} = \frac{8}{1} = 8,\qquad
  L_{22} = \frac{\vA_2^{\!\top}\mat{K}\vA_2}{I_2^2} = \frac{6}{1} = 6,
\]
\[
  L_{12} = \frac{\vA_2^{\!\top}\mat{K}\vA_1}{I_1 I_2} = \frac{2}{1} = 2,
  \qquad L_{21}=L_{12}=2 \ \text{(mirrored, not recomputed).}
\]
The inductance matrix is therefore
\[
  \mat{L} =
  \begin{pmatrix} L_{11} & L_{12}\\ L_{21} & L_{22} \end{pmatrix}
  = \begin{pmatrix} 8 & 2\\ 2 & 6 \end{pmatrix}\si{\nano\henry}.
\]

\textbf{Interpretation.} The diagonal entries are the \emph{self}-inductances:
$L_{11}=\SI{8}{\nano\henry}$ is twice the magnetic energy stored when loop $1$ alone
carries unit current; $L_{22}=\SI{6}{\nano\henry}$ likewise for loop $2$. Both are
positive, as self-inductance must be. The off-diagonal $L_{12}=\SI{2}{\nano\henry}$
is the \emph{mutual} inductance: the flux of loop $1$'s field linked by loop $2$ per
unit current. Here it is positive---the two loops' fields reinforce. Its
\emph{symmetry} $L_{12}=L_{21}$ is reciprocity, and we got it for free: the Gram
computes only the upper triangle and mirrors (\cref{ch:reductions}).

\textbf{The capacitance contrast.} Run \emph{exactly this arithmetic} for the
two-terminal capacitance of \cref{wex:cap2x2} and one thing changes: the weight is
$w=1$ outright (no currents to normalize), the family is potentials $V_i\in\Hone$
not vector potentials $\vA_i\in\Hcurl$, and the operator is the
$\varepsilon$-weighted Poisson stiffness not the $\nu$-weighted curl--curl
stiffness. The off-diagonal also flips \emph{sign} in the capacitance case (the
Maxwell capacitance matrix has negative off-diagonals; mutual inductance need not).
But the reduction \emph{verb} is the same, the symmetry is the same, the
upper-triangle map is the same. Two physical matrices, one Gram.
\end{workedexample}

\subsection{One change at a time: the electrostatic--magnetostatic ledger}

The second worked example is not a computation but an accounting: we lay the two
drivers side by side, line by line, and verify that exactly three lines differ. This
is the chapter's thesis stated as a table.

\begin{workedexample}{Capacitance versus inductance, one change at a time}{ledger}
We tabulate every stage of the two drivers and mark each line ``same'' or
``swap''. Reading down the table, the three swaps stand alone; every other line is
shared verbatim.

\begin{center}\small
\renewcommand{\arraystretch}{1.25}
\begin{tabular}{@{}p{0.20\linewidth}p{0.30\linewidth}p{0.30\linewidth}c@{}}
\toprule
\textbf{stage} & \textbf{electrostatic} & \textbf{magnetostatic} & \\
\midrule
governing PDE & $-\Div(\varepsilon\Grad V)=\rho$ & $\Curl(\nu\Curl\vA)=\vJ$
  & \textcolor{simRust}{\textbf{swap 1}} \\
field unknown & scalar potential $V$ & vector potential $\vA$
  & \textcolor{simRust}{\textbf{swap 1}} \\
post-process & $\vE=-\Grad V$ & $\vB=\Curl\vA$
  & \textcolor{simRust}{\textbf{swap 1}} \\
FE space & $\Hone$ nodal & $\Hcurl$ edge (N\'ed\'elec)
  & \textcolor{simRust}{\textbf{swap 2}} \\
energy operator $\mat{K}$ & $\int\varepsilon\,\Grad\phi_a\!\cdot\!\Grad\phi_b$
  & $\int\nu\,\Curl\boldsymbol\phi_a\!\cdot\!\Curl\boldsymbol\phi_b$
  & \textcolor{simRust}{(swap 2)} \\
\addlinespace
$\mat{K}$ definiteness & SPD (grounded) & SP\emph{S}D (gradient null space)
  & \textcolor{simGray}{setup} \\
solver & CG, fixed op, reused precond. & CG, fixed op, reused precond.
  & \textcolor{simGreen}{\textbf{same}} \\
assemble cadence & once, reused per excitation & once, reused per excitation
  & \textcolor{simGreen}{\textbf{same}} \\
solve cadence & one per terminal & one per source
  & \textcolor{simGreen}{\textbf{same}} \\
reduce verb & \code{gram\_reduce} & \code{gram\_reduce}
  & \textcolor{simGreen}{\textbf{same}} \\
Gram symmetry & $G_{ij}=G_{ji}$, upper triangle & $G_{ij}=G_{ji}$, upper triangle
  & \textcolor{simGreen}{\textbf{same}} \\
\addlinespace
Gram weight $w(i,j)$ & $1$ & $1/(I_iI_j)$
  & \textcolor{simRust}{\textbf{swap 3}} \\
product & capacitance $\mat{C}$ & inductance $\mat{L}$ & \\
\bottomrule
\end{tabular}
\end{center}

\textbf{The reading.} Three rust ``swap'' regions: the PDE-and-field block (swap 1,
which drags post-processing $\vE\to\vB$ along with it), the FE-space line (swap 2,
which determines the energy operator's integrand), and the Gram-weight line (swap
3). The definiteness line is marked grey ``setup'': it is a consequence of swap 2,
but it changes how $\mat{K}$ and $\vJ$ are \emph{prepared} (consistent source or AMS,
\cref{sec:solvecare}), not the solver itself---so the solver line is green
``same.'' Everything in green---the solver, both cadences, the reduce verb, the Gram
symmetry---is shared without modification. The magnetostatic driver is the
electrostatic driver, three swaps over.
\end{workedexample}

\summarysection

Magnetostatics computes an inductance matrix the way electrostatics computes a
capacitance matrix: \emph{the same way}. The driver is the composition
\code{gram\_reduce} $\circ$ \code{solve\_family} $\circ$ \code{fe\_assemble}, and it
differs from the electrostatic driver in exactly three places.

\begin{itemize}[nosep]
\item \textbf{Swap 1---the PDE and field.} Magnetostatics solves the curl--curl
  equation $\Curl(\nu\Curl\vA)=\vJ$ for a magnetic \emph{vector} potential $\vA$,
  not Poisson for a scalar $V$; it post-processes $\vB=\Curl\vA$. Both are
  second-order self-adjoint elliptic equations from adjacent rungs of the de Rham
  complex (\cref{ch:derham}).
\item \textbf{Swap 2---the FE space.} The vector potential lives in $\Hcurl$,
  realized by edge (N\'ed\'elec) elements, because $\vA$'s tangential component is
  continuous while its normal component is free. Nodal-vector elements would impose
  too much continuity and breed spurious modes (\cref{ch:bc}); edge elements sit in
  the exact discrete de Rham complex and avoid them. The curl--curl operator
  inherits a large gradient null space (the gauge fields), making it only positive
  semi-definite; the solve handles this through a consistent ($\Div\vJ=0$) source or
  AMS preconditioning (\cref{ch:smoothers})---a change of \emph{setup}, not of
  solver.
\item \textbf{Swap 3---the Gram weight.} Inductance is the same \code{gram\_reduce}
  as capacitance with the current-normalized weight $w=1/(I_iI_j)$ in place of
  $w=1$ (\cref{ch:reductions}). One verb, two physical quantities; the diagonal is
  twice the magnetic energy per unit current ($L_{ii}=2U_m/I^2$), the off-diagonal
  is the (signed, symmetric) mutual inductance.
\end{itemize}

The shared remainder---assemble once, solve per source with the same CG, reduce by a
symmetric Gram---is the deeper lesson: a simulator earns its generality by making
distinct physical analyses \emph{the same machine over different vocabulary}. The
next chapter (\cref{ch:eigenmodes}) reuses the very same curl--curl operator on the
very same $\Hcurl$ space, now as an eigenvalue problem, and the gradient null space
of \cref{sec:nullspace} returns as the subspace the eigensolver must project out.

\furtherreading

The curl--curl equation, the magnetic vector potential, and edge-element
discretization of magnetostatics are developed in Jin's computational
electromagnetics text~\citep{jin2014} and, with full finite-element-theoretic rigor
(N\'ed\'elec spaces, the discrete de Rham complex, and convergence), in
Monk~\citep{monk2003}. The physical content---vector potential, magnetic energy,
self- and mutual inductance---is the magnetostatics of any intermediate
electromagnetism course; Griffiths~\citep{griffiths2017} is the standard reference,
and connects the inductance matrix to the circuit parameters an engineer expects.
The auxiliary-space Maxwell (AMS) preconditioner for the curl--curl gradient null
space is due to Hiptmair~\citep{hiptmair1998}; \cref{ch:smoothers} treats it in the
multigrid context.

\exercisessection

\begin{enumerate}[label=\textbf{40.\arabic*}]

\item \textbf{The three swaps from memory.} Without looking back, write the
electrostatic driver as a composition \code{reduce} $\circ$ \code{solve\_family}
$\circ$ \code{fe\_assemble} and then derive the magnetostatic driver from it by
applying the three swaps. For each swap, name (i) the box it lives in and (ii) one
sentence of why the physics forces it.

\item \textbf{Curl--curl is self-adjoint.} Show that the bilinear form
$a(\vA,\vF)=\int_\dom\nu\,(\Curl\vA)\cdot(\Curl\vF)\dV$ is symmetric,
$a(\vA,\vF)=a(\vF,\vA)$, and hence that the assembled stiffness $\mat{K}$ of
\cref{eq:stiff-curlcurl} is a symmetric matrix. Why does this symmetry matter for
(a) the conjugate-gradient solver and (b) the Gram reduction's reciprocity
$L_{ij}=L_{ji}$?

\item \textbf{The null space is gradients.} Let $\psi_h$ be any discrete nodal field
and $\mat{G}$ the discrete gradient (\cref{ch:derham}, \cref{def:discrete-grad}).
Show $\mat{K}\,(\mat{G}\psi_h)=\mathbf 0$ directly from
\cref{eq:stiff-curlcurl}, using $\Curl\Grad=0$. Conclude that $\mat{K}$ is only
positive semi-definite, with null space $\supseteq\operatorname{range}(\mat{G})$.

\item \textbf{Why a consistent source rescues CG.} The curl--curl null space is the
range of the discrete gradient. Argue that a divergence-free source $\vJ$
($\Div\vJ=0$) is orthogonal to every gradient field, hence has no component along
the null space, hence conjugate gradients applied to $\mat{K}\vA=\vJ$ stays in the
positive-definite complement and converges. (Hint: integrate $\int\vJ\cdot\Grad\psi$
by parts.)

\item \textbf{The current normalization.} Suppose loop $1$ is driven at
$I_1=\SI{2}{\ampere}$ and loop $2$ at $I_2=\SI{3}{\ampere}$ (not unit currents), and
the solve returns $\vA_1^{\!\top}\mat{K}\vA_1=48$,
$\vA_2^{\!\top}\mat{K}\vA_2=54$, $\vA_2^{\!\top}\mat{K}\vA_1=18$ (energy-form units).
Compute the inductance matrix $\mat{L}$ using the weight $w(i,j)=1/(I_iI_j)$. Verify
$L_{12}=L_{21}$ and that the values are the same as if you had first rescaled each
field to unit current.

\item \textbf{Self-inductance is positive; mutual need not be.} Explain, from
$L_{ii}=\vA_i^{\!\top}\mat{K}\vA_i/I_i^2$ and the positive-semidefiniteness of
$\mat{K}$, why $L_{ii}\ge 0$ always. Then construct a physical scenario (two loops,
relative orientation) in which $L_{12}<0$, and contrast with the capacitance matrix,
whose off-diagonals are always negative (\cref{ch:reductions}, \cref{wex:cap2x2}).

\item \textbf{Nodal vector would break it.} Describe, in two or three sentences, what
goes wrong if one discretizes $\vA$ with three copies of nodal $\Hone$ elements
instead of $\Hcurl$ edge elements. Identify the property of the discrete de Rham
complex (\cref{ch:derham}) that fails and the symptom it produces in the spectrum
(\cref{ch:bc}).

\item \textbf{The eigenmode connection (forward to \cref{ch:eigenmodes}).} The
eigenmode driver assembles the \emph{same} curl--curl $\mat{K}$ on the \emph{same}
$\Hcurl$ space, but solves $\mat{K}\vx=\lambda\mat{M}\vx$ instead of
$\mat{K}\vA=\vJ$. Predict: which eigenvalues does the gradient null space of
\cref{sec:nullspace} produce, and why must the eigensolver project them out? (You
are reconstructing the spurious-zero-mode discussion of \cref{ch:bc} from the
null-space structure alone.)

\end{enumerate}
